\scanpage{17}\mspage{16}
une relation d'ordre ; en d'autres termes

(i) $\forall x\in M$, dans l'ens ${\unsure{\mathcal{F}|\mathscr{F}}}^{x}$ des $F\in\mathcal{F}$ telles que $\unsure{X|x}\in F$, il y a un plus petit \unsure{élément}, soit $F_{x}$ \unsure{[|(}NB ce sera l'ens des $Y\in M$ tels que $Y\leq x$)

(ii) L'application $X\unsure{\mathrel{\text{---}}|\longrightarrow} F_{X}$ est injective, i.e. si $X,Y\in M$, $X\ne Y$, alors $\exists F\in\mathcal{F}$ telle que $X\in F,\ Y\notin F$, ou l'inverse $Y\in F,\ X\notin F$.

Dès lors, les $F\in\mathcal{F}$ \deleted{\unsure{sont des}} \inserted{forment un ens} de parties fermées de $M$, \inserted{contenant les $M_{\leq x}$, $x\in M$\unsure{!?|!|?}} \deleted{\unsure{Se donner un} $\mathcal{F}\subset\unsure{\mathcal{P}|\mathrm{P}}(M)$ \unsure{satisfais.|satisfaisant} (i) (ii), \unsure{revient au même que de se donner}} \inserted{De plus, on suppose} $\varnothing$ (partie vide de $M$) $\in\mathcal{F}$ (ce qui implique que $\mathcal{F}$ n'est \uline{jamais} admissible !).

Dans GF $\overline{\underline{\mathrm{IV}}}$ (Analysis Situs, \unsure{1ère|prem.} mouture) j'ai explicité en long et en large la relation entre la donnée de $(\mathcal{F},\leq)$, \unsure{satisfais.|satisfais} à des conditions qui y étaient \unsure{appelées|appellées} $\underline{\mathrm{C1}}$, $\underline{\mathrm{C2}}$ $\underline{\mathrm{C3}}$ (\unsure{pages|page} \deleted{\unsure{52|53}}\unsure{53|52} \deleted{\unsure{56|54}}\unsure{54|56}, 58),
% slanted
\margin{J'appellerai ces conditions \unsure{plutôt} $\underline{\mathrm{At}}\,1$, $\underline{\mathrm{At}}\,2$, $\underline{\mathrm{At}}\,3$}
\deleted{\unsure{que}} et un ens ordonné $M$, muni d'un ens de parties fermées $\mathcal{F}$, contenant les $M_{\leq \deleted{\unsure{x|y}}x}$, la partie vide, \inserted{et} \unsure{avec} \unsure{tout|toute} $F$ ses parties fermées (pour l'ordre de $M$), plus la condition : si $F,G\in\unsure{\mathcal{J}|\mathcal{F}}$, tels que $\forall x\in F,\ \unsure{y|\forall y}\in G$ on ait $\unsure{\overset{\scriptscriptstyle\cdots}{\scriptscriptstyle v}|\overset{\scriptscriptstyle\cdots}{\scriptscriptstyle\vee}}\,M_{\leq x}\cup M_{\leq y}\in\mathcal{F}$, alors $F\cup G\in\mathcal{F}$. (critère de compatibilité).
